\documentclass[11pt]{article}

\usepackage[margin=0.68in]{geometry}
\usepackage{amsmath,amssymb}
\usepackage{enumitem}
\usepackage{xcolor}
\usepackage{tcolorbox}
\usepackage{fancyhdr}
\usepackage{lastpage}
\usepackage{hyperref}
\usepackage{microtype}
\usepackage{array}
\usepackage{booktabs}

\definecolor{uwpurple}{HTML}{4B2E83}
\definecolor{uwgold}{HTML}{B7A57A}
\definecolor{softpurple}{HTML}{F4F1FA}
\definecolor{softgold}{HTML}{FBF7EA}
\definecolor{deepgold}{HTML}{8A6F2A}
\definecolor{softgray}{HTML}{F6F6F6}

\tcbuselibrary{breakable,skins}

\hypersetup{
    colorlinks=true,
    linkcolor=uwpurple,
    urlcolor=uwpurple
}

\pagestyle{fancy}
\fancyhf{}
\lhead{\textbf{MATH 394: Probability I}}
\rhead{\textbf{Final Exam Preparation Guide}}
\cfoot{\thepage\ of \pageref{LastPage}}
\setlength{\headheight}{14pt}

\setlength{\parindent}{0pt}
\setlength{\parskip}{0.45em}

\newcommand{\course}{MATH 394: Probability I}
\newcommand{\instructor}{Arman Jahangiri}
\newcommand{\department}{Department of Mathematics}
\newcommand{\university}{University of Washington}

\newtcolorbox{topicbox}[1][]{
    breakable,
    colback=softpurple,
    colframe=uwpurple,
    boxrule=0.8pt,
    arc=2mm,
    left=2.5mm,
    right=2.5mm,
    top=1.5mm,
    bottom=1.5mm,
    title={#1},
    coltitle=white,
    colbacktitle=uwpurple,
    fonttitle=\bfseries
}

\newtcolorbox{toolbox}[1][]{
    breakable,
    colback=softgold,
    colframe=deepgold,
    boxrule=0.8pt,
    arc=2mm,
    left=2.5mm,
    right=2.5mm,
    top=1.5mm,
    bottom=1.5mm,
    title={#1},
    coltitle=uwpurple,
    colbacktitle=softgold,
    fonttitle=\bfseries
}

\newtcolorbox{warningbox}[1][]{
    breakable,
    colback=white,
    colframe=uwgold,
    boxrule=0.8pt,
    arc=2mm,
    left=2.5mm,
    right=2.5mm,
    top=1.5mm,
    bottom=1.5mm,
    title={#1},
    coltitle=uwpurple,
    colbacktitle=softpurple,
    fonttitle=\bfseries
}

\begin{document}

% ============================================================
% TITLE PAGE
% ============================================================

\begin{titlepage}
\centering
\vspace*{0.6cm}

{\Large \textbf{\university}}\\
{\large \department}

\vspace{0.8cm}

{\Huge \color{uwpurple}\textbf{Final Exam Preparation Guide}}\\[0.22cm]
{\Large \textbf{\course}}\\[0.18cm]
{\large Summer 2026 \quad | \quad Instructor: \instructor}

\vspace{0.8cm}

\begin{tcolorbox}[
    width=0.88\textwidth,
    colback=softpurple,
    colframe=uwpurple,
    boxrule=0.9pt,
    arc=3mm
]
\centering
\renewcommand{\arraystretch}{1.3}
\begin{tabular}{>{\bfseries}p{0.30\textwidth}p{0.48\textwidth}}
Final exam date: & August 21, 2026 \\
Answering time: & 60 minutes \\
Number of problems: & 7 \\
Total points: & 60 \\
\end{tabular}
\end{tcolorbox}

\vspace{0.6cm}

\begin{warningbox}[Purpose of this guide]
This document tells you  what topic each problem
tests, what mathematical tools you should know how to use, and what
background calculations you should be comfortable performing.

 
\end{warningbox}

\vfill

{\large\color{uwpurple}\textbf{The final emphasizes conceptual recognition,
correct setup, and efficient calculation.}}

\end{titlepage}

% ============================================================
% OVERVIEW
% ============================================================

\section*{1. Exam}

\renewcommand{\arraystretch}{1.35}
\begin{center}
\begin{tabular}{>{\centering\arraybackslash}p{0.10\textwidth}
                >{\centering\arraybackslash}p{0.11\textwidth}
                p{0.55\textwidth}
                >{\centering\arraybackslash}p{0.12\textwidth}}
\toprule
\textbf{Problem} & \textbf{Points} & \textbf{Main topic} & \textbf{Approx. time} \\
\midrule
1 & 16 & Joint distributions, covariance, independence, MGF & 13 min \\
2 & 8 & Multivariate transformation and Jacobian & 9 min \\
3 & 7 & Jensen, symmetry, Cauchy--Schwarz & 6 min \\
4 & 7 & Chebyshev inequality and sample-size control & 7 min \\
5 & 8 & Convolution of continuous random variables & 7 min \\
6 & 8 & Pairwise versus mutual independence & 8 min \\
7 & 6 & WLLN, SLLN, and CLT (multiple choice) & 4 min \\
\midrule
\textbf{Total} & \textbf{60} & & \textbf{54 min} \\
\bottomrule
\end{tabular}
\end{center}

\begin{warningbox}[Advicew]
The approximate solving times total about \(54\) minutes. This leaves several
minutes for reading, checking algebra, fixing notation, and returning to a
part that you initially skipped.
\end{warningbox}

% ============================================================
% PROBLEM 1
% ============================================================

\newpage
\section*{2. Problem-by-Problem Preparation}

\begin{topicbox}[Problem 1 \hfill 16 points \hfill Joint distributions, covariance, and MGF]

\textbf{What this problem covers}

\begin{itemize}[leftmargin=*]
    \item Expectations from marginal densities.
    \item Expectations from a joint density.
    \item Covariance:
    \[
    \operatorname{Cov}(X,Y)
    =
    E(XY)-E(X)E(Y).
    \]
    \item Covariance of affine transformations.
    \item Determining whether two random variables are independent.
    \item Calculating a moment-generating function from a density.
    \item Recovering the expectation from the derivative of the MGF.
\end{itemize}

\end{topicbox}

\begin{toolbox}[Tools you must know for Problem 1]

You should be able to use
\[
E(X)
=
\int_{-\infty}^{\infty}x f_X(x)\,dx,
\]
and
\[
E(XY)
=
\int_{-\infty}^{\infty}
\int_{-\infty}^{\infty}
xy f_{X,Y}(x,y)\,dx\,dy.
\]

You should know
\[
\operatorname{Cov}(X,Y)
=
E(XY)-E(X)E(Y),
\]
and the affine covariance rule
\[
\operatorname{Cov}(aX+b,cY+d)
=
ac\,\operatorname{Cov}(X,Y).
\]
 

You should therefore also remember
\[
\operatorname{Var}(X)
=
E(X^2)-[E(X)]^2.
\]

For independence, you should know that jointly continuous \(X,Y\) are
independent if and only if
\[
f_{X,Y}(x,y)=f_X(x)f_Y(y)
\]
on the support.

You should also remember the implication
\[
X\perp Y
\quad\Longrightarrow\quad
\operatorname{Cov}(X,Y)=0
\]
whenever the relevant moments exist.

The converse is \textbf{not} true in general.
\end{toolbox}

\begin{toolbox}[MGF tools]

The definition is
\[
M_X(t)
=
E(e^{tX})
=
\int_{-\infty}^{\infty}
e^{tx}f_X(x)\,dx.
\]

You should know
\[
M_X'(0)=E(X),
\]
and, more generally,
\[
M_X^{(k)}(0)=E(X^k).
\]

Be careful to determine the values of \(t\) for which the MGF exists.
\end{toolbox}

\begin{toolbox}[Calculus background needed]

You should be comfortable with improper integrals, integration by parts, and substitution.
 
\end{toolbox}

\begin{warningbox}[Common mistakes]
\begin{itemize}[leftmargin=*]
    \item Forgetting that \(E(XY)\) requires the \textbf{joint} density.
    \item Confusing zero covariance with independence.
    \item Forgetting that constants disappear inside covariance.
    \item Forgetting to check the domain of \(t\) for an MGF.
\end{itemize}
\end{warningbox}

% ============================================================
% PROBLEM 2
% ============================================================

\newpage

\begin{topicbox}[Problem 2 \hfill 8 points \hfill Transformation of two random variables]

\textbf{What this problem covers}


Jacobian method for finding the transformation of random variables.

\end{topicbox}

\begin{toolbox}[The complete Jacobian procedure]

Suppose
\[
Z=g_1(X,Y),
\qquad
W=g_2(X,Y).
\]

You should follow these steps.

\begin{enumerate}[leftmargin=*]
    \item Write the transformation:
    \[
    z=g_1(x,y),
    \qquad
    w=g_2(x,y).
    \]

    \item Solve for the original variables:
    \[
    x=x(z,w),
    \qquad
    y=y(z,w).
    \]

    \item Determine the new support in the \((z,w)\)-coordinates.

    \item Compute the inverse Jacobian:
    \[
    J
    =
    \frac{\partial(x,y)}{\partial(z,w)}.
    \]

    \item Take the absolute value:
    \[
    |J|.
    \]

    \item Use
    \[
    f_{Z,W}(z,w)
    =
    f_{X,Y}(x(z,w),y(z,w))
    \left|
    \frac{\partial(x,y)}{\partial(z,w)}
    \right|.
    \]

    \item To obtain \(f_Z(z)\), integrate out \(W\):
    \[
    f_Z(z)
    =
    \int f_{Z,W}(z,w)\,dw.
    \]
\end{enumerate}

\end{toolbox}
 

% ============================================================
% PROBLEM 3
% ============================================================

\newpage
\begin{topicbox}[Problem 3 \hfill 7 points \hfill Comparing expectations and probabilities]

This problem is conceptual. You need to recognize which theorem or
inequality applies.

The main tools you should be able to recognize are:

\begin{itemize}[leftmargin=*]
    \item Jensen's inequality;
    \item Markov's inequality;
    \item Chebyshev's inequality;
    \item Cauchy--Schwarz inequality;
    \item symmetry of i.i.d.\ random variables.
\end{itemize}

\end{topicbox}

\begin{toolbox}[Jensen's inequality]

If \(g\) is convex,
\[
g(E(X))
\leq
E[g(X)].
\]

If \(g\) is concave, the direction reverses:
\[
E[g(X)]
\leq
g(E(X)).
\]

You should therefore know how to determine convexity or concavity using
\[
g''(x).
\]

In particular,
\[
g(x)=\log x
\]
is concave for \(x>0\), since
\[
g''(x)=-\frac1{x^2}<0.
\]

\end{toolbox}

\begin{toolbox}[Markov's inequality]

If \(X\) is a nonnegative random variable and \(a>0\), then
\[
P(X\geq a)
\leq
\frac{E(X)}{a}.
\]

More generally, Markov's inequality may be applied to any
\textbf{nonnegative function of a random variable}. For example,
\[
P(X^2\geq a)
\leq
\frac{E(X^2)}{a}.
\]

Before applying Markov's inequality, always check that the random variable
to which you apply it is nonnegative.

\end{toolbox}

\begin{toolbox}[Chebyshev's inequality]

If \(X\) has finite mean and variance, then for every \(a>0\),
\[
P\left(
|X-E(X)|\geq a
\right)
\leq
\frac{\operatorname{Var}(X)}{a^2}.
\]

Equivalently,
\[
P\left(
|X-E(X)|<a
\right)
\geq
1-
\frac{\operatorname{Var}(X)}{a^2}.
\]

You should recognize that Chebyshev's inequality controls deviations
\textbf{from the mean}, whereas Markov's inequality directly controls the
upper tail of a nonnegative random variable.

Also remember that Chebyshev's inequality can be obtained by applying
Markov's inequality to
\[
(X-E(X))^2.
\]

\end{toolbox}
% ============================================================
% PROBLEM 4
% ============================================================

\newpage

\begin{topicbox}[Problem 4 \hfill 7 points \hfill Chebyshev inequality and sample size]

\textbf{What this problem is about}

\begin{itemize}[leftmargin=*]
    \item Mean and variance of a sample average.
    \item Chebyshev's inequality.
    \item Converting an upper bound on an error probability into a lower
    bound on a success probability.
    \item Choosing an efficient sample size.
\end{itemize}

\end{topicbox}
 

\begin{toolbox}[Chebyshev's inequality]

For any random variable with finite variance,
\[
P(|X-E(X)|\geq a)
\leq
\frac{\operatorname{Var}(X)}{a^2}.
\]

You should be comfortable switching between
\[
P(|X-\mu|\geq a)
\]
and
\[
P(|X-\mu|<a)
\]
using complements.
\end{toolbox}
 
 

% ============================================================
% PROBLEM 5
% ============================================================

\newpage

\begin{topicbox}[Problem 5 \hfill 8 points \hfill Convolution]

\textbf{What this problem covers}

Finding the distribution of the sum of two independent variables through the convolution technique.

\end{topicbox}

\begin{toolbox}[Convolution formula]

For independent continuous random variables,
\[
f_{X+Y}(t)
=
\int_{-\infty}^{\infty}
f_X(x)f_Y(t-x)\,dx.
\]

The most important part is often not the integration itself, but determining
where both densities are simultaneously nonzero.
\end{toolbox}


% ============================================================
% PROBLEM 6
% ============================================================

\newpage

\begin{topicbox}[Problem 6 \hfill 8 points \hfill Pairwise and mutual independence]

\textbf{What this problem covers}

\begin{itemize}[leftmargin=*]
    \item Marginalizing a three-dimensional joint density.
    \item Obtaining pairwise joint densities.
    \item Testing pairwise independence.
    \item Testing mutual independence.
    \item Understanding the logical relationship between pairwise and mutual
    independence.
\end{itemize}

\end{topicbox}

\begin{toolbox}[Obtaining pairwise densities]

From
\[
f_{X,Y,Z}(x,y,z),
\]
you should know that
\[
f_{X,Y}(x,y)
=
\int_{-\infty}^{\infty}
f_{X,Y,Z}(x,y,z)\,dz.
\]

Similarly,
\[
f_{X,Z}(x,z)
=
\int
f_{X,Y,Z}(x,y,z)\,dy,
\]
and
\[
f_{Y,Z}(y,z)
=
\int
f_{X,Y,Z}(x,y,z)\,dx.
\]
\end{toolbox}

\begin{toolbox}[Pairwise independence]

The variables \(X,Y,Z\) are pairwise independent if all three conditions
hold:
\[
f_{X,Y}(x,y)=f_X(x)f_Y(y),
\]
\[
f_{X,Z}(x,z)=f_X(x)f_Z(z),
\]
and
\[
f_{Y,Z}(y,z)=f_Y(y)f_Z(z).
\]

If even one pair fails, then the variables are not pairwise independent.
\end{toolbox}

\begin{toolbox}[Mutual independence]

For jointly continuous \(X,Y,Z\), mutual independence, \textbf{in addition to pairwise independence}, requires
\[
f_{X,Y,Z}(x,y,z)
=
f_X(x)f_Y(y)f_Z(z)
\]
on the support.

Remember:
\[
\text{mutual independence}
\quad\Longrightarrow\quad
\text{pairwise independence}.
\]

The converse is not true in general.
\end{toolbox}


% ============================================================
% PROBLEM 7
% ============================================================

\newpage

\begin{topicbox}[Problem 7 \hfill 6 points \hfill WLLN, SLLN, and CLT]

This is a theorem-recognition problem.

You need to distinguish:

\[
\text{WLLN}
\qquad
\text{vs.}
\qquad
\text{SLLN}
\qquad
\text{vs.}
\qquad
\text{CLT}.
\]

You need to know both the \textbf{assumptions} and the \textbf{conclusion} of
each theorem.
\end{topicbox}

\begin{toolbox}[Weak Law of Large Numbers]

Under the assumptions used in class, let
\[
X_1,X_2,\dots
\]
be i.i.d.\ with
\[
E(X_i)=\mu
\]
and finite variance.

Then
\[
\overline X_n\xrightarrow{P}\mu.
\]

Equivalently, for every \(\varepsilon>0\),
\[
\lim_{n\to\infty}
P\left(
|\overline X_n-\mu|>\varepsilon
\right)
=
0.
\]

\textbf{Key phrase:}
\[
\boxed{\text{convergence in probability}}
\]
\end{toolbox}

\begin{toolbox}[Strong Law of Large Numbers]

Under the same assumptions used in class,
\[
\overline X_n\xrightarrow{a.s.}\mu.
\]

Equivalently,
\[
P\left(
\lim_{n\to\infty}\overline X_n=\mu
\right)
=
1.
\]

\textbf{Key phrase:}
\[
\boxed{\text{almost-sure convergence}}
\]

Also remember:
\[
\text{almost-sure convergence}
\quad\Longrightarrow\quad
\text{convergence in probability}.
\]

Thus the SLLN is stronger than the WLLN.
\end{toolbox}

\begin{toolbox}[Central Limit Theorem]

Let
\[
X_1,X_2,\dots
\]
be i.i.d.\ with
\[
E(X_i)=\mu,
\qquad
\operatorname{Var}(X_i)=\sigma^2<\infty.
\]

Then
\[
\frac{\overline X_n-\mu}{\sigma/\sqrt n}
\xrightarrow{d}
N(0,1).
\]

Equivalently, for every \(z\in\mathbb R\),
\[
\lim_{n\to\infty}
P\left(
\frac{\overline X_n-\mu}{\sigma/\sqrt n}
\leq z
\right)
=
\Phi(z).
\]

\textbf{Key phrase:}
\[
\boxed{\text{convergence in distribution}}
\]
\end{toolbox}

\begin{warningbox}[Do not confuse the three conclusions]

\[
\begin{array}{c|c|c}
\textbf{Theorem}
&
\textbf{What converges?}
&
\textbf{Type of convergence}
\\
\hline
\text{WLLN}
&
\overline X_n\to\mu
&
\text{in probability}
\\[1mm]
\text{SLLN}
&
\overline X_n\to\mu
&
\text{almost surely}
\\[1mm]
\text{CLT}
&
\dfrac{\overline X_n-\mu}{\sigma/\sqrt n}
\to N(0,1)
&
\text{in distribution}
\end{array}
\]

The WLLN/SLLN13 tells you \textbf{where the sample mean goes}.

The CLT tells you the \textbf{shape of the fluctuations around the mean}.
\end{warningbox}

% ============================================================
% MASTER TOOLKIT
% ============================================================
 

\end{document}